\section{Problemas encontrados y comentarios}

\subsection{Incidencias y soluciones}

\begin{table}[H]
  \centering
  \begin{tabularx}{\textwidth}{>{\raggedright\arraybackslash\bfseries}p{0.27\textwidth}XX}
    \toprule
    Incidencia & Causa & Solución aplicada \\
    \midrule
    Cierre de DBeaver con \texttt{SIGBUS} & El fallo ocurrió en la biblioteca
    nativa \texttt{libgtk-3.so.0} durante el inicio con Java 25. & Se utilizó
    OpenJDK 21 y se inició DBeaver con \texttt{GDK\_BACKEND=x11}. \\
    \addlinespace
    Puerto 5432 ocupado & Ya existía un contenedor \texttt{postgres} publicado
    en ese puerto. & Se mantuvo el contenedor existente y se publicó
    \texttt{postgres-practica01} en el puerto 5433. \\
    \addlinespace
    Fallo de autenticación en DBeaver & La conexión conservaba una contraseña
    distinta de la generada para el nuevo contenedor. & Se actualizó la
    credencial guardada en la conexión y se probó nuevamente. \\
    \bottomrule
  \end{tabularx}
  \caption{Problemas observados durante la instalación y conexión.}
\end{table}

\evidencia{05_evidencia.png}{Estado de los contenedores después de resolver el conflicto del puerto y validación de PostgreSQL.}

\subsection{Comentarios finales}

La práctica permitió comprobar el funcionamiento de Docker en Debian y
desplegar PostgreSQL sin instalar el servidor directamente en el sistema
anfitrión. El volumen independiente mantiene los datos del contenedor y la
publicación en la interfaz local evita exponer el servicio a otros equipos de
la red.

La validación se realizó en tres niveles: ejecución de \texttt{hello-world},
consulta directa con \texttt{psql} y conexión mediante DBeaver. Además, el
conflicto del puerto mostró la diferencia entre el puerto publicado en el
anfitrión (5433) y el puerto interno del servidor (5432).

\textbf{Resultado.} Docker quedó habilitado y activo; el contenedor
\texttt{postgres-practica01} permaneció funcionando con PostgreSQL 18.6; y
DBeaver se conectó correctamente mediante el puerto local 5433.
